%% [AI DRAFT] Borrador generado para discusión (Modo C, etapa 3).
%% Citas verificadas contra el PDF. Sin marcadores [VERIFY] pendientes.

\section{CONTEXTO}

Esta sección delimita el marco conceptual y técnico que sustenta la arquitectura desarrollada. 
Para ello, se fundamenta el entorno económico de aplicación, se define la naturaleza de los atributos extraíbles, 
se contrasta el método de lectura frente a las técnicas de recolección automatizada, y se establecen 
las reglas de interpretación para el manejo de evidencia insuficiente.

\subsection{Comercio digital y confianza en el Perú}

El comercio electrónico peruano opera en un entorno donde la adopción del pago digital no está acompañada por un nivel equivalente de confianza del consumidor. 
Un análisis comparado de sistemas de pago en economías emergentes reporta que la falta de confianza en las instituciones es una barrera clave, provocando que
solo el 57\,\% de los adultos utilice pagos digitales, frente al 64\,\% global \cite{ref26}.
Esta desconfianza genera un impacto directo en el mercado legítimo, ya que la intención y continuidad de compra en canales digitales dependen fuertemente del riesgo percibido por el usuario
\cite{ref21} \cite{ref6}. 
Por lo tanto, la incapacidad de verificar a un vendedor desconocido antes del pago no solo afecta a las víctimas de fraude, 
sino que impone una fricción comercial sobre los vendedores honestos.

\subsection{Características extraíbles del lado del cliente}

Para efectos de esta arquitectura, se define como \emph{característica extraíble del lado del cliente} a todo atributo de una publicación que cumpla tres condiciones simultáneas: es visible para el comprador antes del pago, reside en el documento ya entregado al navegador, y no requiere acceso a interfaces internas de la plataforma. 
Atributos como la puntuación del producto, el número de reseñas y el nombre de dominio cumplen estrictamente con estos criterios.

Esta distinción separa el presente trabajo de la línea dominante en detección de fraude, 
la cual se apoya en telemetría interna (registros de transacciones, huellas de dispositivo, grafos de
interacción) completamente inaccesible para el comprador \cite{ref28}. Por su parte,  el comprador dispone de un
conjunto de características restringido, por lo que la pregunta central que este trabajo aborda es si dichos datos accedibles bastan para 
producir una advertencia útil.

\subsection{Lectura del DOM y recolección automatizada}

La operación que realiza la extensión requiere una caracterización precisa, con consecuencias
técnicas y no solo terminológicas. Una revisión sistemática de noventa y un estudios define el
\textit{web scraping} y el \textit{web crawling} como el proceso automatizado de extracción de
datos en el que las herramientas de rastreo navegan a través de múltiples páginas
\cite{ref24}. La extensión propuesta no realiza ninguna de las dos operaciones: no navega,
porque es el usuario quien decide qué ficha abrir; no recorre múltiples páginas, porque opera
sobre el documento que el navegador ya renderizó; y no emite peticiones adicionales hacia la
plataforma. Su operación consiste en interpretar, en beneficio de quien la ejecuta, contenido
que ya le había sido entregado.

Esta caracterización acota también las limitaciones del enfoque. La misma revisión advierte
que las herramientas apoyadas en selectores predefinidos son susceptibles a cambios de
maquetación, y que entre los desafíos del área figuran los riesgos legales \cite{ref24}. La
primera observación se declara como amenaza a la validez en la Sección~V-C; la segunda es la
razón para no retransmitir contenido de la plataforma más allá de las características que el
veredicto requiere.

\subsection{Protección de pago declarada}

Las plataformas de comercio de mayor tamaño publican esquemas mediante los cuales retienen el
pago hasta la conformidad del comprador o garantizan su devolución ante incumplimiento. Esa
declaración es un compromiso público y verificable en la propia ficha, y constituye por ello
una característica legítima.

Su ausencia, en cambio, no es informativa en el mismo sentido: un vendedor independiente
puede ofrecer condiciones equivalentes sin anunciarlas, o anunciarlas donde la extensión no
alcanza a leer. Tratar el silencio como evidencia de inseguridad convertiría una limitación
del observador en una acusación sobre el observado, tanto más cuanto que las contramedidas
disponibles se despliegan y se financian del lado del comercio \cite{ref30} y su ausencia
visible correlaciona con el tamaño del vendedor antes que con su honestidad.
